\section{Additional posterior analyses}
\label{app:posterior-analysis}

\subsection{Posterior learning across recursive backbones}
\label{sec:cross-backbone}

\Cref{tab:maze-cross-backbone} compares posterior methods with TRM and FPRM backbones on the complete Maze-Hard test set.
Every stochastic method uses ten candidates ($K=10$) and its native selector.

\begin{table}[t]
  \caption{\textbf{FB obtains the highest selected accuracy and Pass@10
  with both backbones on Maze-Hard.} Results are percentages on the complete
  1,000-puzzle test set. Bold and underlined values mark the best and
  second-best results within each backbone.}
  \label{tab:maze-cross-backbone}
  \centering
  \small
  \setlength{\tabcolsep}{22pt}
  \begin{adjustbox}{max width=1.0\linewidth}
  \begin{tabular}{lcc}
    \toprule
    Method & Acc. $\uparrow$ & Pass@10 $\uparrow$ \\
    \midrule
    \multicolumn{3}{l}{\textit{TRM backbone}} \\
    Deterministic            & 87.00 & 87.00 \\
    PTRM                     & \underline{87.70} & 88.40 \\
    W-PTRM                   & 87.50 & \underline{89.80} \\
    \rowcolor{gaerowhighlight}
    \textbf{FB (Ours)}       & \textbf{88.00} & \textbf{91.00} \\
    \addlinespace
    \multicolumn{3}{l}{\textit{FPRM backbone}} \\
    Deterministic            & \underline{86.90} & -- \\
    W-FPRM                   & \underline{86.90} & 88.70 \\
    IVON                     & 84.40 & \underline{94.00} \\
    \rowcolor{gaerowhighlight}
    \textbf{FB (Ours)}       & \textbf{87.70} & \textbf{94.80} \\
    \bottomrule
  \end{tabular}
  \end{adjustbox}
\end{table}

FB improves both metrics under the two tested backbones.
With TRM, it raises selected accuracy by 0.30 percentage points over PTRM and Pass@10 by 1.20 points over W-PTRM.
With FPRM, FB improves selected accuracy from 84.40\% to 87.70\% and Pass@10 from 94.00\% to 94.80\% over IVON.
The paired 95\% bootstrap intervals are [1.60, 5.10] for selected accuracy and [0.00, 1.60] for Pass@10.
The two FPRM posterior rows use the same initialization, training data, update budget, and evaluation noise.

\subsection{Posterior construction under a fixed backbone}
\label{sec:posterior-construction}

\Cref{tab:posterior-construction} fixes the FPRM backbone and Maze-Hard protocol while varying how each method constructs its parameter distribution.
Diagonal Laplace fits a local Gaussian from curvature at a point estimate~\citep{daxberger2021laplace}, while SWAG estimates a Gaussian from an SGD trajectory~\citep{maddox2019swag}.

\begin{table}[h]
  \caption{\textbf{The comparison fixes FPRM and Maze-Hard while varying
  the posterior construction.}}
  \label{tab:posterior-construction}
  \centering
  \small
  % \begin{adjustbox}{width=1.0\linewidth}
  \begin{tabular}{lcc}
    \toprule
    Method &  Acc. $\uparrow$ & Pass@10 $\uparrow$ \\
    \midrule
    FPRM                    & \underline{86.90} & -- \\
    W-FPRM                  & \underline{86.90} & 88.70 \\
    Diagonal Laplace        & \pendingcell{Diagonal Laplace Acc.} & \pendingcell{Diagonal Laplace Pass@10} \\
    SWAG                    & \pendingcell{SWAG Acc.} & \pendingcell{SWAG Pass@10} \\
    IVON                    & 84.40 & \underline{94.00} \\
    \rowcolor{gaerowhighlight}
    \textbf{FB (Ours)}      & \textbf{87.70} & \textbf{94.80}\\
    \bottomrule
  \end{tabular}
  % \end{adjustbox}
\end{table}

Among posterior results, FB improves selected accuracy by 3.30 percentage points and Pass@10 by 0.80 points over IVON.
\resultpending{Compare FB with Diagonal Laplace and SWAG after their matched Maze-Hard results and fit times are available.}

\paragraph{Candidate diversity.}
\placeholder{Report the unique-solution ratio and candidate disagreement for
each posterior method. Pair each diversity statistic with Pass@$K$ so that
output variation is not interpreted as useful coverage by itself.}

